\section{Introduction}

A cellular automaton is an infinite array of identical finite machines, the cells, each updating its state in
step with all the others by one fixed local rule that reads only the cell and its neighbours. The question of
universality asks whether such a fixed local rule, with a fixed finite set of states, can carry out any
computation at all. In the Euclidean plane the question was settled long ago, most famously by the Game of
Life, where gliders are launched, steered, and collided to build logic. The hyperbolic plane is a harder and
richer arena. It admits infinitely many regular tilings rather than the three of the Euclidean plane, and its
straight lines diverge, so the glider trick fails. Two gliders aimed at each other in hyperbolic space drift
apart and never meet.

Margenstern resolved the difficulty with a different idea, the railway \cite{stewart1994}. Instead of
colliding moving patterns, one lays down a fixed network of tracks and lets a single locomotive run on it
forever. The locomotive never collides with anything and never needs to be synchronised with a second
locomotive. Where it passes through a switch it may change the switch, and the settings of all the switches,
frozen in place, are the memory of the computation. A suitable network of tracks and switches reproduces a
register machine, and a register machine with two counters is already universal \cite{minsky1967}. The
railway therefore gives universality with nothing to choreograph, which is exactly what hyperbolic space
rewards.

On this foundation Margenstern, in part with collaborators, built a sequence of concrete automata
\cite{margenstern2008vol2}. The pentagrid, the tiling $\{5,4\}$ by right-angled pentagons four to a vertex,
received a rotation-invariant weakly universal automaton, with five states \cite{margenstern2014penta5},
then three \cite{margenstern2015penta3}, then two \cite{margenstern2015penta2}. The ternary heptagrid
$\{7,3\}$ received its own automaton \cite{margenstern2009hepta}. Hyperbolic three-dimensional space, tiled by
right-angled dodecahedra in the dodecagrid $\{5,3,4\}$, received a five-state automaton
\cite{margenstern2021dodeca5} and an outer totalistic one with four states \cite{margenstern2021dodeca4}. Each
of these is a genuine achievement. Each is also a separate achievement. Every tiling was treated on its own,
with its own neighbourhood, its own table of rules, and its own machine-checked verification, because the
local picture, which neighbour is which and how a rotation permutes them, is different in each tiling and even
differs from cell to cell within one tiling.

\subsection{The gap}

The state of the art is therefore a family of points, one universal automaton per tiling, obtained one at a
time. The regular tilings of the hyperbolic plane are infinite in number, $\{p,q\}$ for every pair with
$1/p + 1/q < 1/2$, and the buildable honeycombs of hyperbolic three-space and four-space add more. A natural
question is whether the points can be replaced by a single object. Is there one cellular automaton, one fixed
rule, that is universal on every regular tiling, rather than a different automaton for each?

The obstruction the per-tiling work overcame is the absence of a global frame. A cell cannot name a fixed
compass direction, and the number and arrangement of its neighbours change with the tiling. Margenstern met
this with rotation invariance and with milestones, small decorations that mark the orientation of a track
locally. The price of rotation invariance was paid anew in every tiling.

\subsection{The contribution}

We give one cellular automaton, defined by rules that are local in the cell-adjacency graph and never refer to
a particular tiling, that is universal on every regular hyperbolic tiling at once, and we make it \textit{strongly}
universal, so that a computation can be set up from a finite configuration rather than an infinite one. The
automaton realizes the railway. Its moving part is a four-symbol alphabet, empty space, clear track, a head,
and a trailing cell, together with a single bit carried by each switch. The locomotive advances correctly with
no global frame for a simple reason made precise in Section~\ref{sec:automaton}. The cell immediately behind the
head is the trailing cell of the locomotive, never clear track, so the only clear track cell adjacent to a head
lies ahead of it, and that is the cell the head moves into. Direction is read off the local pattern, not off
the plane.

The step from weak to strong universality is the heart of the paper. A weakly universal automaton is allowed an
infinite initial configuration, and the unbounded registers of a register machine are usually supplied that
way, as an infinite ultimately-periodic track. We instead let the registers \textit{build themselves}. The locomotive,
on overflowing a register, lays a new piece of track by a local construction rule and so adds a new digit. A
register that began as a single cell grows without bound as the computation needs it, exactly as a growing mesh
adds a new ring at its open edge each step. The program is finite, the registers self-extend, so the whole
machine starts from a finite seed, which is strong universality. Section~\ref{sec:strong} gives this in full,
with the complete state set and transition tables.

The cost we pay for covering every tiling with one rule is that our switches are oriented. A switch cell
carries, as part of the initial configuration, the labels of its trunk and its two branches. Because the labels
are local data, not a global frame, the automaton stays tiling-independent, but it is not rotation invariant.
We are explicit about this tradeoff in Section~\ref{sec:comparison}. Margenstern bought rotation invariance and
small state counts at the price of solving each tiling separately. We buy one strongly universal construction
for all tilings at the price of oriented switches.

The plan of the paper is as follows. Section~\ref{sec:railway} recalls the railway and its three switches.
Section~\ref{sec:setting} fixes the combinatorial setting, a tiling seen as its cell-adjacency graph.
Section~\ref{sec:automaton} defines the automaton, gives the full state set and transition tables, and proves
its motion and switching correct. Section~\ref{sec:strong} is the central section, the self-building registers
and the strong universality theorem. Section~\ref{sec:universality} assembles the register machine.
Section~\ref{sec:verification} reports the verification on the actual cell graphs of seven tilings.
Section~\ref{sec:weak} describes the weaker variant. Section~\ref{sec:comparison} compares the construction with
the tiling-specific automata. Section~\ref{sec:conclusion} concludes. Appendices~\ref{app:rules} and~\ref{app:tables} collect the complete transition table and worked traces. The construction and all of its checks live in the accompanying code
\cite{pollard2026vibe}.
